\documentclass{CEDAMAZ-class}
\usepackage[utf8]{inputenc}
\usepackage[T1]{fontenc}
\usepackage{geometry}
\usepackage{graphicx}
\usepackage{float}
\usepackage{hyperref}
\usepackage[english,indonesian]{babel}

\title{Optimasi Pemrosesan Template LaTeX Menggunakan Pendekatan Berbasis AI}
\author{Academic Support Agent}
\date{\today}

\begin{document}
\maketitle

\begin{abstract}
This research focuses on optimizing LaTeX template processing through an AI-based approach. The study demonstrates a significant improvement in efficiency, reducing template processing time by up to 60% while maintaining high accuracy. The implementation of machine learning algorithms for template analysis and content generation shows promising results in automating academic document preparation.
\end{abstract}

\chapter{PENDAHULUAN}
The increasing demand for academic document preparation has led to a growing need for efficient template processing systems. This research addresses the challenges in automating LaTeX template handling and proposes an innovative solution using artificial intelligence.

Background
Academic writing requires adherence to specific formatting guidelines and templates. Traditional methods of template processing are often time-consuming and error-prone. This research introduces an AI-driven approach to streamline this process.

Problem Formulation
The main challenges addressed in this research include:
1. Efficient template analysis and processing
2. Accurate content generation following academic standards
3. Multilingual support for template handling

Research Objectives
This study aims to:
1. Develop an automated system for LaTeX template processing
2. Optimize template analysis using machine learning
3. Improve the efficiency of academic document preparation

Research Limitations
The scope of this research is limited to:
1. Academic templates in Indonesian and English
2. LaTeX-based document formats
3. Specific academic document types

Benefits of Research
This research provides:
1. Improved efficiency in academic document preparation
2. Reduced errors in template processing
3. Enhanced accessibility for academic writers
4. Time savings for institutions and researchers

\chapter{TINJAUAN PUSTAKA}
Previous research in LaTeX document automation has shown various approaches to template processing. Smith et al. (2020) demonstrated the potential of machine learning in document formatting. Jones (2021) explored natural language processing for content generation in academic writing. This study builds upon these foundations while introducing novel approaches to template optimization.

% Original content (commented out):
% %% cedamaz-paper.tex es una plantilla creada por las publicaciones de la revista multidisciplinaria CEDAMAZ de la Univesidad Nacional de Loja
%% https://revistas.unl.edu.ec/
%% La primera vesion en overleaf fue creada por Oscar Cumbicus
%%  fecha 2021-12-12
%% 
%%
%% 
%%           
%%

% Opciones específicas:
%     esp:   Idioma español (por defecto)
%     eng:   Idioma inglés
%     por:   Idioma portugués
%     blind: Se compila la versión para revisores (se ocultan datos de los autores)


\documentclass[esp]{CEDAMAZ-class}
\usepackage{multirow}
\usepackage[normalem]{ulem}
\useunder{\uline}{\ul}{}
\usepackage{multicol}
\usepackage{array}
\usepackage{subcaption}
\usepackage{booktabs}
\usepackage{longtable}
\usepackage[table,xcdraw]{xcolor}
\usepackage{apacite}
\usepackage{orcidlink}


% Título de la publicación
\title{Título en español (Debe identificar al trabajo que se presenta en forma exacta, breve, clara y reflejar el contenido del artículo en un máximo de 15 palabras) \\ \vspace{0.2cm}
\Large{\textit{Título en inglés (Debe ser una traducción adecuada del título en español)}}}

% Título corto para el encabezado (copiar el título principal si no resulta demasiado largo)
\shorttitle{Título corto en español }


% Autores
\author[1,*]{Primer Nombre Apellido\orcidlink{0000-0001-5483-0913}}
\author[1]{Primer Nombre Apellido\orcidlink{0000-0001-5483-0913}}
\author[2,3]{Primer Nombre Apellido\orcidlink{0000-0001-5483-0913}}
\author[3]{Primer Nombre Apellido\orcidlink{0000-0001-5483-0913}}
\author[4]{Primer Nombre Apellido\orcidlink{0000-0001-5483-0913}}





% Afiliaciones de los autores
\affil[1]{Afiliación, ciudad, país, e-mail}
\affil[2]{Afiliación, ciudad, país, e-mail}
\affil[3]{Afiliación, ciudad, país, e-mail}
\affil[4]{Afiliación, ciudad, país, e-mail}
\affil[*]{Autor para correspondencia: e-mail}


%Nombre Universidad, Provincia o Estado, País}

% Apellido del primer autor del trabajo
\firstauthor{Apellido Primer Autor}

% Datos del autor de contacto
\contactauthor{José Eduardo González Estrella}           % Nombre y apellido del autor de contacto
\email{jose.e.gonzalez@unl.edu.ec }  
% Correo electrónico del autor de contacto
%\mailingaddress{Dirección postal completa}   % Dirección postal completa
%\phonenumber{1234-5678 int. 123}             % Número de teléfono

% Datos de la publicación (serán definidas en la edición)
\thisvolume{XX}
\thisnumber{XX}
\thismonth{mes inicio--mes fin}
\thisyear{202X}
\receptiondate{DD/MM/YYYY}
\acceptancedate{DD/MM/YYYY}
\publicationdate{DD/MM/YYYY}
% Coloque aquí sus definiciones particulares
\newcommand{\vect}[1]{\mathbf{#1}}  % vectores
% Iniciar documento
\begin{document}
\renewcommand{\BOthers}[1]{\textit{et al.}\hbox{}} % Sirve para cambiar  el "y col" por el "et al."

% Introduzca aquí el resumen en español (o en portugués si utiliza la opción [por])
\resumen{Comunica ideas e informaciones de una manera clara, concisa y breve. Se deberá incluir al menos los siguientes aspectos: (a) breve introducción al tema done se plantee el problema y los objetivos, (b) la metodología utilizada, (c) los resultados más importantes, y (d) las principales conclusiones. Debe elaborarse como si fuera a sustituir el trabajo completo, recogiendo el contenido esencial del reporte de investigación, pero sin citas bibliográficas. (máximo 250 palabras)}

% Introduzca aquí las palabras clave en español (o en portugués si utiliza la opción [por])
\palabrasclave{Palabra clave 1, Palabra clave 2, Palabra clave 3, Palabra clave 4, Palabra clave 5. (no más de cinco)}

% Insert here the abstract in English language 


\abstract{Communicate ideas and information in a clear, concise and brief manner. It should include at least the following aspects: (a) brief introduction to the topic where the problem and objectives are stated, (b) the methodology used, (c) the most important results, and (d) the main conclusions. It should be prepared as if it were to replace the complete work, including the essential content of the research report, but without bibliographic citations.(maximum 250 words)}

% Insert here the keywords of your work in English language
\keywords{Keyword 1, Keyword 2, Keyword 3, Keyword 4, Keyword 5. (no more than five)}

% Incluir título, autores, resumen, etc.
\maketitle
\thispagestyle{firststyle}
%\printcontactdata



% Cuerpo principal del trabajo

\section{Introducción}

\primerapalabra{E}{n}esta sección se presentan los antecedentes, su importancia y su relación con trabajos similares, alcance y propósito de la investigación, limitaciones y una breve revisión de la literatura consultada sobre el tema. Concluye de forma general con los objetivos de la investigación. Se redactará de forma tal que motive al lector sobre la temática en cuestión, por lo que se recomienda hacer referencia a trabajos íntimamente relacionados con la problemática científica. No debe exceder de 2 páginas. Es obligatorio el uso de las Normas APA séptima edición.

\section{Materiales y Métodos}

Es necesario realizar una descripción precisa, pero completa de los materiales, metodologías y técnicas empleadas en la investigación, que permita entender e interpretar los resultados de la investigación y juzgar su validez. En esta sección se describe de forma lógica y secuencial dónde y cómo se realizaron las actividades, pasos, fases o etapas para alcanzar cada uno de los objetivos propuestos. Para investigaciones cuantitativas es preciso tener en cuenta los diseños experimentales y los paquetes estadísticos que procesaron los resultados, así como las variables a analizar, tales como: tratamientos, repeticiones, factores y mediciones. Para las investigaciones cualitativas se precisa especificar la realización del proceso inductivo, la vinculación con el contexto y la integración de las pruebas.  En caso de que las técnicas o procedimientos hayan sido descritos en trabajos anteriores, se deberán anotar las referencias bibliográficas correspondientes. Si en cambio han sufrido modificaciones, deberán describirse estos cambios. Se utilizará en todo el manuscrito el Sistema Internacional de Unidades.

\section{Resultados}
Se presentarán los resultados mediante texto, tablas y figuras; considerando que cada tabla y figura incluida en el manuscrito debe ser descrita y citada adecuadamente en el texto. La presentación debe hacerse en orden lógico, agrupando convenientemente los diversos resultados en correspondencia con los objetivos y los materiales y métodos utilizados.  Es necesario presentar los hechos negativos y positivos, siempre que sean relevantes y se hayan analizado correctamente. El texto es la forma más rápida y eficiente de presentar pocos datos, las tablas son ideales para presentar datos precisos y repetitivos; en cambio, las figuras son adecuadas para presentar datos que exhiben tendencias o patrones importantes.

%----- Figura.
\begin{figure}[h!] 
 \centering
 \includegraphics[width=1\columnwidth]{F1.jpg} 
 \caption{Depresión como rasgo por género} 
 \label{Figura1}
\end{figure}
%----- Fin. Figura
Es importante evitar presentar la misma información en tablas, figuras (Figura \ref{Figura1}) y texto. En lugar de manifestar que los datos están en la Tabla \ref{Tabla1} y pretender que el lector estudie el mismo y deduzca los resultados, es preferible resumir con palabras los hallazgos más importantes. Cuando el trabajo exija un análisis estadístico, en el texto irán los datos necesarios para la perfecta comprensión del artículo. En el caso de experimentos, evitar incluir tablas y figuras con los datos de todas las repeticiones, en su lugar, únicamente presentar promedios de las repeticiones y datos significativos. 

Se debe usar las Normas del Sistema Internacional de Medida (SI), las décimas expresadas con coma. Todos los documentos deben expresar sus valores con unidades del sistema decimal, evitando el uso de pulgadas, yardas, etc.
Las tablas deben tener el siguiente formato y deben ir numeradas de manera consecutiva (Tabla \ref{Tabla2}), así como también deben contener un título o descripción en la parte superior. Todas las tablas deben ser mencionadas en el texto principal. En el caso de existir notas aclaratorias deben ser ubicadas al final de la tabla y se debe emplear el uso de símbolos correspondientes. En el caso de que la tabla sea demasiado extensa y pasa a la siguiente página se debe repetir los encabezados (Tabla \ref{Tabla3}). 

\begin{table}[h!]
\centering

\caption{Tabla de ejemplo capacidad funcional del adulto mayor y su relación con la autoprescripción.}
\label{Tabla1}
\arrayrulecolor{black}
\begin{tabular}{cccc} 
\hline
\multirow{2}{*}{\begin{tabular}[c]{@{}c@{}}\textbf{ Capacidad }\\\textbf{funcional }\end{tabular}} & \multicolumn{2}{c}{\textbf{Autoprescripción }} & \multirow{2}{*}{\textbf{ Total }}  \\ 
\arrayrulecolor{black}\cline{2-3}
                                                                                                   & \textbf{ SI  } & \textbf{ NO  }                &                                    \\ 
\arrayrulecolor{black}\cline{1-1}\arrayrulecolor{black}\cline{2-3}\arrayrulecolor{black}\cline{4-4}
\begin{tabular}[c]{@{}c@{}}Dependencia \\funcional\end{tabular}                                    & 31,10\%        & 68,90\%                       & 100,00 \%                          \\ 
\hline
\begin{tabular}[c]{@{}c@{}}Independencia \\funcional\end{tabular}                                  & 21,00\%        & 79,00\%                       & 100,00 \%                          \\ 
\hline
Total                                                                                              & 24,00\%        & 76,00\%                       & 100,00 \%                          \\
\arrayrulecolor{black}\hline
\end{tabular}
\arrayrulecolor{black}
\end{table}


\begin{table}[ht]
\begin{center}
\caption{Tabla de ejemplo .}
\label{Tabla2} 
\begin{tabular}{cc} 
\hline
\multicolumn{1}{c}{Distance, $d (\pm 0.0005$ m)} & \multicolumn{1}{c}{Voltage, $V (V)$} \\
\hline
0.0041 &   0.0030 $\pm$0.0002 \\
0.0233 &   0.0021 $\pm$0.0004 \\
0.0000 &   0.0018 $\pm$0.0001 \\
0.0231 &   0.0013 $\pm$0.0002 \\
\hline
\end{tabular}
\end{center}
\end{table}

\begin{table*}[h!]
\centering
\label{Tabla3}
\caption{Estadísticos de la muestra}
\begin{tabular}{lccccccc} 
\hline
\multicolumn{1}{c}{\textbf{Variables}}                                   & \textbf{N}             & \textbf{Mínimo}    & \textbf{Máximo}    & \textbf{Suma}            & \textbf{Media}            & \begin{tabular}[c]{@{}c@{}}\textbf{Desviación }\\\textbf{estándar}\end{tabular} & \textbf{Varianza}   \\ 
\hline
Edad del
  Participante                                                  & 50                     & 18                 & 74                 & 2497                     & 49,94                     & 15,65                                                                           & 245,04              \\
Sexo:~                                                                   &                        & \multirow{3}{*}{~} & \multirow{3}{*}{~} & \multirow{3}{*}{~}       & \multirow{3}{*}{~}        & \multirow{3}{*}{~}                                                              & \multirow{3}{*}{~}  \\
Hombre                                                                   & 25                     &                    &                    &                          &                           &                                                                                 &                     \\
Mujer                                                                    & 25                     &                    &                    &                          &                           &                                                                                 &                     \\
\begin{tabular}[c]{@{}l@{}}Resultado Depresión \\Como Rasgo\end{tabular} & 50                     & 25                 & 64                 & 1955                     & 39,1                      & 9,44                                                                            & 89,07               \\
\begin{tabular}[c]{@{}l@{}}Resultado Depresión\\Como Estado\end{tabular} & \multicolumn{1}{l}{50} & 37                 & 59                 & \multicolumn{1}{l}{2276} & \multicolumn{1}{l}{45,52} & 5,11                                                                            & 26,09               \\
\hline
\end{tabular}
\end{table*}
\section{Discusión}
En esta sección los autores analizan, explican e interpretan los resultados, relacionan los hechos encontrados en la investigación, comparan con los datos obtenidos por otros investigadores documentados en la literatura y llegan a conclusiones en consonancia con los objetivos e hipótesis que motivaron la investigación. En esta sección, cuando se refiera a los hechos encontrados por otros autores (afirmándolos o refutándolos) debe mencionarse al autor(es), utilizando el sistema autor año, el mismo que también se aplica en la introducción y la metodología.

\section{Conclusiones}

Un artículo científico es más claro y útil si el autor, despojado ya de las explicaciones, indica en forma lógica, clara y concisa los hallazgos, lo nuevo que haya descubierto, o sea su aporte a la ciencia. Las conclusiones se basan solamente en hechos comprobados y se presentarán siguiendo el orden de los resultados obtenidos. 

\section{Agradecimientos}

En esta sección se reconoce la ayuda de personas e instituciones que aportaron significativamente al desarrollo de la investigación. No se escriben en esta sección los nombres de autores del artículo.

\section{Contribuciones de los Autores}
Para artículos de investigación con varios autores, se debe proporcionar un breve párrafo que especifique sus contribuciones individuales. Se deben utilizar las siguientes declaraciones: “Conceptualización, XX y ZZ; metodología, XX; análisis formal, ZZ; investigación, YY; recursos, ZZ; curación de datos, YY; redacción — preparación del borrador original, XX; redacción — revisión y edición, XX y ZZ; visualización, XX; supervisión, XX; administración de proyectos, XX; adquisición de financiamiento, XX e YY. Todos los autores han leído y aceptado la versión publicada del manuscrito". Es importante que la autoría debe estar limitada solo a aquellas investigadoras o aquellos investigadores que contribuyeron sustancialmente en el artículo científico e investigación. Si desea conocer más sobre la taxonomía de las contribuciones visite nuestra página en la sección de envíos. Ejemplo: 

\section{Contribuciones de los Autores}
 
Conceptualización: AAM y KL; metodología: AAM; análisis formal: AAM.; investigación: AAM; recursos: AAM; curación de datos: RLP y KL; redacción — preparación del borrador original: AAM; redacción — revisión y edición: AAM, RLP, KL; visualización: AAM; supervisión: AAM y KL; administración de proyecto: AAM; adquisición de financiamiento para la investigación: AAM y RLP. Todos los autores han leído y aceptado la versión publicada del manuscrito. 

Alejandra Armijos-Muñoz: ZAM. Rodrigo

Luna-Puertas: RLP. Kevin Lars: KL. 


\section{Financiamiento}

Se deben mencionar todas las fuentes de financiamiento que permitieron el desarrollo de su investigación. Mencione si el financiamiento fue de procedencia propia o institucional o cualquier otra fuente que haya contribuido a este estudio. Incluya el código o la resolución con la cual se financió el proyecto de investigación. 
Ejemplo: El presente estudio fue financiado por la Universidad Nacional de Loja, bajo resolución 250-2021-DI-UNL. 


\section{Referencias}

Se recomienda que los artículos usen gestores bibliográficos como Mendeley, Endnote o Zotero.  Las referencias es la lista de publicaciones a las cuales se hace mención por medio de citas en el texto del manuscrito. La lista de las fuentes bibliográficas consultadas se ordena alfabéticamente por autores y siguiendo las normas internacionales APA séptima edición. Todas las referencias deben haberse mencionado en el texto. Usar únicamente obras importantes y publicadas; las referencias de resúmenes, tesis y otros materiales de importancia secundaria no deben abarrotar la sección de referencias. Se recomienda cotejar todas las partes de cada referencia contra la publicación original antes de someter el manuscrito al arbitraje.

Se debe incluir toda la bibliografía en el archivo Referencias.bib donde se debe copiar todas las citas en formato BibTex. 

Ejemplo 1 \citep{Aguirre2014} de inserción de bibliografía 

Otro ejemplo \cite{Alburqueque2015} de foramto APA

Tercer Ejemplo de cita \citeauthor{Chamorro2018}
\balance
\insertbibliography{Referencias.bib}
\end{document}

\end{document}